\scanpage{103}

\uline{Remarques} Si $\unsure{E|\mathcal E}$ est loc. libre de rg $n+2$, alors on trouve des conditions équivalentes $(\overline{\mathrm A})$, \deleted{\unsure{$\overline{\mathrm B}$|$(\mathrm B)$}} et $(\overline{\mathrm B})$, en renforçant (A), (B), remplaçant \deleted{\unsure{(8)}}(9) de (A) et (12) de (B) respectivement par

\[(s_i)_{0\leq i\leq n+1}\text{ \unsure{forme|forment} une base de }E\tag*{(9$'$)}\]

\[(a_i)_{-1\leq i\leq n}\text{ \unsure{forme|forment} une base de }E.\tag*{(12$'$)}\]

\deleted{Quand ces conditions} C'est aussi équivalent à la conjonction de (A), ou (B), avec la condition

\[E_{n+1}=E\tag*{(26)}\]

\deleted{Soit alors $a'_{0}=a'_{-1}$} \deleted{Si ces conditions sont satisfaites, \unsure{complétons} \unsure{Soit} Soit \unsure{$a'_j$|$a'_{-1}$} \inserted{$s'_0=a'_{-1}$}, une forme lin. sur $E$ telle que}

\[\deleted{\operatorname{Ker}a'_{-1}=k(a_0,\ldots,a_n)}\tag*{\deleted{(27)}}\]

\deleted{elle est donc déterminée, modulo scalaire $\in k^*$, par la seule donnée des $\tau_i$ $(0\leq i\leq n)$, du}

Idem pour cor. 1 [conditions $(\overline{\mathrm A}'),(\overline{\mathrm B}')$, équivalent à (A$'$) ou (B$'$) $+{}$(26)]

\uline{Proposition} \unsure{Soit|Soient} $E$ un $k$-module loc. libre de rg $n+2$, $\tau_i$ $(0\leq i\leq n)$ des pseudo-réflexions dans $E$,

\[\tau_i=\mathrm{id}+a'_i\otimes a_i\qquad 0\leq\unsure{i|0}\leq n,\tag*{(27)}\]

\deleted{satisfaisant (28) $\tau_i\tau_j=\tau_j\tau_i$ dans $\underline{\mathrm P}(E)$.}

Pour qu'il existe \inserted{localement} \unsure{un|des} $s_0\in E$ tel que les conditions $(\overline{\mathrm A})(\overline{\mathrm B})$ du Th. 1 \inserted{(remarque ci-dessus)} soient satisfaites (conditions \inserted{commutation et} de ``non dég''), il f et s.

% archive sheet number: 102
